% Part: counterfactuals
% Chapter: introduction
% Section: paradoxes-material

\documentclass[../../../include/open-logic-section]{subfiles}

\begin{document}

\olfileid{cnt}{int}{par}

\olsection{ભૌતિક શરતી વિધાનના વિરોધાભાસો}

અંગ્રેજીના ``જો \dots તો \dots'' વાક્યોનું પ્રતીકીકરણ કરવા
$!A \lif !B$ વાપરવાની ટીકા સૌપ્રથમ કરનારાઓમાં સી.~આઈ. લૂઈસ હતા.
તેઓ વ્હાઇટહેડ અને રસેલના \emph{પ્રિન્સિપિયા મેથેમેટિકા}માં ભૌતિક
શરતી વિધાનના ઉપયોગની ટીકા કરતા હતા; ત્યાં $\lif$ને ``ફલિત કરે છે''
એમ વાંચવામાં આવતું હતું. લૂઈસે યોગ્ય રીતે વાંધો ઉઠાવ્યો કે જો
$\lif$નો અર્થ ``ફલિત કરે છે'' હોય, તો કોઈ પણ અસત્ય વિધાન~$p$,
$p$માંથી~$q$ ફલિત થાય છે એવું ફલિત કરે છે, કેમ કે $p$ અસત્ય હોય
ત્યારે $p \lif (p \lif q)$ સત્ય છે. એ જ રીતે કોઈ પણ સત્ય વિધાન~$q$,
$p$માંથી~$q$ ફલિત થાય છે એવું ફલિત કરે છે, કેમ કે $q$ સત્ય હોય
ત્યારે $q \lif (p \lif q)$ સત્ય છે.

તર્કશાસ્ત્રીઓ જાણે છે કે \emph{ફલિતતા}, એટલે કે તાર્કિક ફલિતતા,
સંયોજક નથી; તે !!{formula}s અથવા વિધાનો વચ્ચેનો સંબંધ છે. તેથી
ગૂંચવણ ટાળવા $\lif$ને ``ફલિત કરે છે'' એમ વાંચવું ન જોઈએ.\footnote{$\lif$ને
  ``ફલિત કરે છે'' એમ વાંચવાનો પ્રયોગ ગણિતશાસ્ત્રીઓ અને કમ્પ્યુટર
  વૈજ્ઞાનિકોમાં હજુ વ્યાપક છે; જોકે લૂઈસે દર્શાવેલી ગૂંચવણો ટાળવા
  તત્ત્વચિંતકો તેને ``માત્ર ત્યારે જ'' એમ વાંચે છે.} આ ભેદ રાખીએ
તો લૂઈસની ખાસ ચિંતા ઊભી થતી નથી: $p$ને અસત્ય અંગ્રેજી વાક્યનું
પ્રતીક માનીએ તોપણ $p$, $q$ને ``ફલિત'' કરતું નથી. $p \Entails q$
છે કે નહિ તે નક્કી કરવા \emph{બધાં} !!{valuation}s ધ્યાનમાં લેવાં પડે
છે; $p$ વડે હકીકતમાં અસત્ય હોય એવું વાક્ય પ્રતીકિત કરીએ તોપણ
$p \Entails/ q$ છે.

છતાં લૂઈસના નીચેના વાક્ય જેવાં ``જો \dots તો\dots'' વાક્યોમાં
કંઈક વિચિત્ર લાગે છે:
\begin{quote}
જો ચંદ્ર લીલા પનીરનો બનેલો હોય, તો $2+2=4$.
\end{quote}
તેમજ નીચેના અનુમાનોમાં પણ:
\begin{quote}
  ચંદ્ર લીલા પનીરનો બનેલો નથી. તેથી, જો ચંદ્ર લીલા પનીરનો બનેલો
  હોય, તો $2+2=4$.

  $2+2 = 4$. તેથી, જો ચંદ્ર લીલા
  પનીરનો બનેલો હોય, તો $2+2=4$.
\end{quote}
પરંતુ ``જો \dots તો \dots''નો અર્થ માત્ર $\lif$ જ હોય, તો આ વાક્ય
કોઈ મુશ્કેલી વિના સત્ય અને આ અનુમાનો કોઈ મુશ્કેલી વિના માન્ય ગણાય.

બીજી ચિંતા પુનરુક્તિ $(!A \lif !B) \lor (!B \lif
!A)$માંથી ઊભી થાય છે. તે એવું સૂચવે છે કે અખબારમાંથી કોઈ પણ બે
નિવેદક વાક્યો~$S$ અને $T$ યાદૃચ્છિક રીતે લઈએ, તો ``જો $S$ તો
$T$, અથવા જો $T$ તો~$S$'' વાક્ય સત્ય હોવું જોઈએ.

\end{document}
